\documentclass[11pt,a4paper]{article}

\usepackage[spanish,es-noshorthands]{babel}
\usepackage[utf8]{inputenc}
\usepackage[T1]{fontenc}
\usepackage{amsmath}
\usepackage[margin=2.5cm]{geometry}
\usepackage{longtable}
\usepackage{booktabs}
\usepackage{array}
\usepackage{xcolor}
\usepackage{hyperref}
\usepackage{titlesec}
\usepackage{enumitem}
\usepackage{fancyhdr}
\usepackage{parskip}

\definecolor{pantone186}{RGB}{200,16,46}

\hypersetup{
  colorlinks=true,
  linkcolor=pantone186,
  urlcolor=pantone186,
  citecolor=pantone186,
  pdftitle={Modelamiento de carga a la base de datos - usuario comun},
  pdfauthor={Equipo Ticketera Nestjs}
}

\titleformat{\section}{\normalfont\Large\bfseries\color{pantone186}}{\thesection}{1em}{}
\titleformat{\subsection}{\normalfont\large\bfseries}{\thesubsection}{1em}{}

\pagestyle{fancy}
\setlength{\headheight}{14pt}
\fancyhf{}
\renewcommand{\headrulewidth}{0.4pt}
\fancyhead[L]{Modelamiento de carga a la base de datos}
\fancyhead[R]{\thepage}
\fancyfoot[C]{Documento interno — confidencial}

\newcolumntype{L}[1]{>{\raggedright\arraybackslash}p{#1}}

\title{\color{pantone186}\textbf{Modelamiento de carga a la base de datos}\\[0.3em]
\Large Según el perfil de un usuario común (rol Agente)}
\author{Proyecto Ticketera Nestjs}
\date{22 de septiembre de 2026 \\ Versión 1.0}

\begin{document}

\maketitle
\thispagestyle{empty}
\newpage
\tableofcontents
\newpage

% =====================================================================
\section{Objetivo y alcance}

Este documento estima la carga que un usuario típico del sistema genera
sobre PostgreSQL, y la contrasta contra los límites de infraestructura ya
fijados en el proyecto (pool de conexiones, \emph{timeouts}, tamaño de lote
del \emph{scheduler} de SLA). El propósito es identificar si la
configuración actual —pensada y medida, no adivinada, según
\texttt{CLAUDE.md}— soporta el volumen de uso esperado, y en qué punto deja
de hacerlo.

No es una medición contra tráfico real: el proyecto no tiene aún usuarios
de producción, solo el \emph{seed} de prueba (5 usuarios, ver
\S\ref{sec:supuestos}). Es un modelo estimativo, con supuestos explícitos,
pensado para poder recalcularse cuando existan métricas reales.

% =====================================================================
\section{Perfil de usuario modelado}

Se eligió el rol \textbf{AGENTE} como "usuario común" del sistema, por
sobre REQUESTER y SUPERVISOR/ADMIN, por tres razones:

\begin{itemize}[leftmargin=1.2cm]
  \item Es quien más interactúa con la base durante una jornada completa:
    un REQUESTER típico solo abre el sistema para crear un ticket o
    revisar su estado —sesiones cortas y esporádicas—, mientras que el
    agente trabaja contra el sistema toda su jornada laboral.
  \item Es, numéricamente, el rol más repetido entre los usuarios activos
    de un sistema de mesa de ayuda: hay muchos más agentes atendiendo que
    supervisores o administradores.
  \item Es el rol que dispara la mayor variedad de operaciones sobre
    \texttt{tickets}: listar, ver detalle, comentar, cambiar estado — el
    conjunto completo de escrituras del dominio, salvo la creación
    (típicamente del REQUESTER o de la ingesta de correo) y la
    administración de usuarios (ADMIN).
\end{itemize}

Al final del documento (\S\ref{sec:otros-perfiles}) se compara brevemente
contra REQUESTER, para dimensionar cuánto pesa cada perfil en el total.

% =====================================================================
\section{Supuestos del modelo}
\label{sec:supuestos}

Ningún supuesto de volumen está documentado en el proyecto (el \emph{seed}
solo trae 2 agentes de prueba), así que se fijan explícitamente aquí, con
un criterio conservador de institución mediana:

\begin{itemize}[leftmargin=1.2cm]
  \item Jornada de \textbf{8 horas} (28\,800 segundos) por agente.
  \item \textbf{20 agentes} concurrentes en horario de mayor actividad —
    cifra asumida, ajustable; todos los cálculos de esta sección escalan
    linealmente con este número.
  \item El agente revisa su panel (\emph{dashboard}) al iniciar turno y
    luego \textbf{cada 45 minutos} aproximadamente (refresco manual o
    vuelta a la pestaña): 10 cargas de \emph{dashboard} por turno.
  \item Además del \emph{dashboard}, hace \textbf{15 búsquedas/listados}
    de tickets sueltas por turno.
  \item Abre el \textbf{detalle de 20 tickets} por turno; en un 30\% de
    esos casos puede reasignar (dispara una consulta adicional a
    \texttt{users-service}).
  \item Escribe \textbf{12 comentarios} y realiza \textbf{8 cambios de
    estado} por turno.
  \item No crea tickets ni administra usuarios (eso corresponde a
    REQUESTER y ADMIN respectivamente).
\end{itemize}

% =====================================================================
\section{De acción de usuario a consultas SQL}

Cada acción del agente se traduce en una o más llamadas RPC (BFF o gateway
→ microservicio), y cada llamada en una o más consultas a PostgreSQL. La
tabla siguiente mapea eso, a partir del código real: los \emph{composers}
del BFF (\texttt{apps/bff-web/src/composers/}) y los casos de uso de
\texttt{tickets-service}.

\begin{longtable}{@{}>{\bfseries}p{3.3cm} L{4.6cm} p{1.6cm} p{1.4cm} p{1.3cm}@{}}
\toprule
\textbf{Acción} & \textbf{Qué dispara} & \textbf{Queries/ evento} & \textbf{Frec./ turno} & \textbf{Total} \\
\midrule
\endhead
Inicio de sesión & \texttt{users-service}: \texttt{SELECT} por email + \texttt{UPDATE} de último acceso & 2 & 1 & 2 \\
Carga de \emph{dashboard} & \texttt{AgentDashboardComposer}: 6 llamadas RPC en paralelo (4× \texttt{tickets.search}, \texttt{tickets.statsByAgent}, \texttt{notifications.findForUser}); cada \texttt{search} pagina (consulta + \texttt{COUNT}) & 10 & 10 & 100 \\
Listado/búsqueda suelta & \texttt{TicketListComposer}: \texttt{tickets.search} + \texttt{COUNT} & 2 & 15 & 30 \\
Detalle de ticket & \texttt{TicketDetailComposer}: \texttt{tickets.findById} (con \emph{joins}) + \texttt{users.find-many} si puede reasignar (30\% de los casos) & 1{,}3 & 20 & 26 \\
Comentario & Lectura de autorización + \texttt{INSERT ticket\_comments} en \texttt{tickets-service}; evento \texttt{ticket.commented} dispara en \texttt{notifications-service} resolución de destinatario + \texttt{INSERT notifications} & 4 & 12 & 48 \\
Cambio de estado & \texttt{SELECT ... FOR UPDATE} (lock pesimista) + \texttt{UPDATE ticket} + \texttt{INSERT ticket\_status\_history}; evento \texttt{ticket.status-changed} dispara 2 queries más en \texttt{notifications-service} & 5 & 8 & 40 \\
\midrule
\textbf{Total por agente / turno} & & & & \textbf{246} \\
\bottomrule
\end{longtable}

Con 246 consultas en 28\,800 segundos, el promedio por agente es de
\textbf{\textasciitilde 0{,}0085 consultas/segundo} — una cifra trivial en
sí misma. El riesgo no está en el promedio: está en la concentración de
esas consultas en ráfagas cortas, que se analiza a continuación.

% =====================================================================
\section{Carga agregada y análisis de concurrencia}

\subsection{Carga promedio del equipo}

Con 20 agentes concurrentes bajo los mismos supuestos:
\[
20 \times 246 = 4\,920 \text{ consultas / turno de 8h} \approx 0{,}17 \text{ consultas/segundo en promedio}
\]

Muy por debajo de cualquier límite de \texttt{DB\_POOL\_MAX=10}: en
régimen estacionario, el pool de 10 conexiones sobra con holgura para este
volumen.

\subsection{El problema real: ráfagas, no promedio}

Dos eventos del sistema concentran consultas en ventanas de segundos, y son
los que efectivamente presionan el pool de 10 conexiones:

\begin{enumerate}[leftmargin=1.2cm]
  \item \textbf{Inicio de turno simultáneo.} Si los 20 agentes inician
    sesión y cargan su \emph{dashboard} dentro de una ventana de 2 minutos
    (escenario típico: turno que arranca a una hora fija), se generan
    $20 \times (2 + 10) = 240$ consultas en 120 segundos, con picos
    instantáneos de hasta $20 \times 6 = 120$ llamadas RPC casi
    simultáneas — cada una necesita una conexión del pool mientras dura su
    consulta. Con \texttt{statement\_timeout=5\,000\,ms}, una consulta
    lenta retiene su conexión hasta 5 segundos; si más de 10 agentes caen
    en el mismo instante, las siguientes esperan en cola por una conexión
    libre.
  \item \textbf{Corrida del \emph{scheduler} de SLA}. Corre cada
    \texttt{EVERY\_5\_MINUTES} y procesa hasta \texttt{BATCH\_SIZE=200}
    tickets por barrido. En un escenario realista (no los 200 completos,
    sino ~15 tickets que efectivamente incumplen SLA en una corrida), esto
    dispara 1 \texttt{SELECT} de barrido + 15 \texttt{UPDATE} + 15
    \texttt{INSERT} de historial + eventos a \texttt{notifications-service}
    (~30 consultas más) — unas \textbf{46 consultas concentradas en menos
    de 1 segundo}, cada 5 minutos, \textbf{independientemente} de cuántos
    agentes estén conectados. Si esta corrida coincide con el pico de
    inicio de turno, ambas ráfagas compiten por el mismo pool de 10
    conexiones.
\end{enumerate}

Este es el mismo tipo de riesgo que \texttt{CLAUDE.md} ya documenta para
\texttt{AssignTicketUseCase}: con \texttt{DB\_POOL\_MAX=10}, diez
operaciones concurrentes que retengan una conexión agotan el pool contra sí
mismas. El modelo confirma que el punto de fricción no es el volumen diario
de un agente —trivial—, sino la coincidencia de ráfagas de corta duración.

\subsection{Confrontación con los límites de infraestructura}

\begin{longtable}{@{}>{\bfseries}p{4.3cm} p{3.6cm} L{4.8cm}@{}}
\toprule
\textbf{Límite configurado} & \textbf{Valor} & \textbf{Lectura frente al modelo} \\
\midrule
\endhead
\texttt{DB\_POOL\_MAX} & 10 conexiones & Suficiente en régimen estacionario (0{,}17 consultas/seg); insuficiente si $>$10 operaciones de ráfaga caen en el mismo instante. \\
\texttt{statement\_timeout} & 5\,000\,ms & Acota cuánto puede retener una conexión una consulta lenta durante una ráfaga; no evita la cola, solo el bloqueo indefinido. \\
\texttt{lock\_timeout} & 3\,000\,ms & Relevante en cambios de estado concurrentes sobre el mismo ticket (\texttt{FOR UPDATE}); no aplica a la carga de lectura del \emph{dashboard}. \\
\texttt{THROTTLE\_LIMIT} & 100 req. / 60\,000\,ms (por instancia de gateway) & A nivel HTTP, 20 agentes generando ráfagas de \emph{dashboard} (6 llamadas c/u) están muy por debajo del límite; el \emph{throttling} no es el cuello de botella para este perfil. \\
\texttt{BATCH\_SIZE} SLA & 200 tickets / corrida, cada 5 min & Fuente de carga \emph{independiente} del número de agentes conectados; debe sumarse a cualquier pico de uso humano al dimensionar el pool. \\
\bottomrule
\end{longtable}

% =====================================================================
\section{Otros perfiles, para contraste}
\label{sec:otros-perfiles}

\begin{itemize}[leftmargin=1.2cm]
  \item \textbf{REQUESTER}: sesión típica de 2-3 acciones (crear ticket,
    revisar estado, comentar) en vez de una jornada completa. Carga
    estimada por sesión: 1 \texttt{INSERT} (crear ticket, con su
    clasificación de urgencia y cálculo de SLA en la misma transacción) +
    1-2 lecturas de detalle $\approx$ 4-6 consultas por sesión, muy por
    debajo de las 246 del agente. Aun con un volumen de REQUESTERs varias
    veces mayor al de agentes, su carga total es comparable o menor,
    porque sus sesiones son mucho más cortas y espaciadas.
  \item \textbf{SUPERVISOR/ADMIN}: carga por sesión más alta que la del
    agente cuando genera reportes o pantallas de \texttt{privacy-dossier}
    (portabilidad ARCO) — esta última puede disparar más de 90 consultas
    en una sola solicitud (hasta 50 páginas de 1\,000 tickets por lote,
    más comentarios y notificaciones), pero es una acción infrecuente
    (no diaria), por lo que no domina el modelo agregado.
\end{itemize}

% =====================================================================
\section{Conclusiones y recomendaciones}

\begin{itemize}[leftmargin=1.2cm]
  \item El volumen \emph{promedio} de consultas por agente (\textasciitilde
    0{,}0085 consultas/seg) es irrelevante frente a la capacidad del pool
    actual; \textbf{no justifica, por sí solo, subir \texttt{DB\_POOL\_MAX}}.
  \item El riesgo real está en la \textbf{coincidencia de ráfagas}: inicio
    de turno simultáneo de varios agentes y la corrida del \emph{scheduler}
    de SLA. Si el equipo crece más allá de los \textasciitilde 20 agentes
    asumidos, o si se agregan más corridas periódicas (schedulers
    adicionales), conviene volver a correr este modelo con el número real
    de agentes concurrentes.
  \item Medir esto en producción es más confiable que seguir extrapolando
    el supuesto de 20 agentes: instrumentar el número de conexiones activas
    del pool durante el inicio de turno real daría la cifra que hoy es
    estimada.
  \item Si la medición real muestra agotamiento del pool durante el pico de
    inicio de turno, la primera palanca es \textbf{escalonar el
    \emph{login}} (no técnica, sino operativa: no todos los agentes
    entran a la vez) antes que subir \texttt{DB\_POOL\_MAX}, porque este
    último ya tiene un costo documentado: valores altos multiplicados por
    réplicas de servicio agotan conexiones de PostgreSQL más rápido de lo
    que se gana en concurrencia (ver \texttt{CLAUDE.md}, sección de base de
    datos).
  \item Vale la pena revisar si el \emph{scheduler} de SLA
    (\texttt{EVERY\_5\_MINUTES}) puede desplazarse fuera de los minutos
    típicos de inicio de turno (ej. \texttt{:02}, \texttt{:07}...) para no
    competir con el pico humano por el mismo pool.
\end{itemize}

\end{document}
